% Draft for an ICRA 2027 submission.
% Replace \method, \datasetnames, and every \TODO before submission.
% Do not turn the qualitative claims below into quantitative claims until the
% corresponding experiment and statistical evidence are present in the paper.

\providecommand{\method}{Terrain-SLAM}
\providecommand{\datasetnames}{\textbf{[DATASET NAMES]}}
\providecommand{\TODO}[1]{\textbf{[TODO: #1]}}

\section{Introduction}
\label{sec:introduction}

Reliable localization is a prerequisite for autonomous operation in large,
outdoor environments. Visual--inertial odometry (VIO) provides locally smooth
motion estimates without relying on external infrastructure, but its error is
unbounded over long trajectories. Global Navigation Satellite System (GNSS)
measurements can anchor this locally accurate estimate in a global frame.
Unfortunately, the conditions in which a global reference is most valuable---
urban canyons, vegetation, steep terrain, and operation near large structures---
are also conditions in which satellite geometry, blockage, and multipath can
make the reported position intermittent, biased, or overconfident. A fusion
system that treats every GNSS fix as an equally trustworthy three-dimensional
anchor can consequently transfer a poor measurement directly into the state
estimate.

The vertical channel is particularly challenging. Vision and inertial sensing
constrain short-term relative motion well, yet small attitude and inertial-bias
errors can accumulate into substantial height drift. Conversely, satellite
altitude is typically less accurate than horizontal position and may exhibit
large local errors. This asymmetry becomes important on trajectories with
significant elevation change: assuming approximately constant altitude is no
longer defensible, while following noisy GNSS altitude can corrupt an otherwise
smooth visual--inertial trajectory.

Digital elevation models (DEMs) provide a complementary, globally referenced
description of terrain height. Unlike a constant-height prior, a DEM retains
the spatial variation of the terrain; unlike an instantaneous GNSS altitude
measurement, it is unaffected by the current satellite visibility. A DEM does
not, however, directly give platform altitude. Its usefulness depends on a
global horizontal position at which to query the map, an assumed or measured
sensor height above the terrain, consistent vertical datums, and a treatment of
map resolution and terrain/model uncertainty. These dependencies motivate
integrating DEM information as a probabilistic constraint rather than replacing
the complete pose estimate with a terrain lookup.

In this work, we present \method{}, a terrain-aware extension of the
keyframe-based visual--inertial SLAM system OKVIS2-X~\cite{boche2025okvis2x}.
The estimator augments the existing visual, inertial, loop-closure, and GNSS
factors with a one-dimensional DEM height residual. At a globally referenced
horizontal location, the DEM height and the sensor's nominal height above
ground define an expected vertical position; the residual is weighted by a
configurable height uncertainty and includes the sensor-to-body lever arm. This
formulation deliberately constrains only elevation, leaving horizontal motion
to visual--inertial estimation and accepted GNSS observations. The
implementation also supports rejecting implausible GNSS observations and
bounding recovery after disagreement or dropout, so that unreliable global
measurements need not induce abrupt, unrestricted corrections.

We study the resulting system on \datasetnames{} trajectories spanning
\TODO{state the altitude range, distance, environment types, and GNSS-degradation
protocol}. We compare against default GNSS-enabled OKVIS2-X, MINS,
and VINS-Fusion. The evaluation is designed to separate overall position
accuracy from the paper's central question by reporting absolute trajectory
error, vertical error, horizontal error, relative pose error, and robustness
under \TODO{natural and/or synthetically controlled} GNSS degradation. We do
not assume that a terrain prior is uniformly beneficial: ablations isolate the
DEM factor, GNSS screening/recovery, and their combination, and results are
stratified by terrain slope and GNSS quality.

The contributions of this paper are:
\begin{itemize}
  \item a terrain-aware factor-graph formulation that introduces DEM elevation
        as an uncertainty-weighted, one-dimensional constraint in a
        keyframe-based visual--inertial SLAM estimator;
  \item a robust global-aiding pipeline that combines selective GNSS use with
        bounded recovery while preserving locally smooth visual--inertial
        motion; and
  \item an evaluation against OKVIS2-X with its default GNSS fusion, MINS, and
        VINS-Fusion on trajectories with substantial altitude variation and
        degraded GNSS, including component ablations and separate vertical and
        horizontal error analysis.
\end{itemize}

\paragraph{Drafting note.}
The contribution statements above are intentionally non-quantitative. Replace
them with measured claims only after fixing the dataset splits, GNSS-degradation
protocol, estimator configurations, and number of repeated runs. In particular,
``more robust'' should be tied to a declared failure criterion and ``more
accurate'' to confidence intervals or paired sequence-level statistics.

\section{Related Work}
\label{sec:related_work}

\subsection{Visual--inertial odometry and SLAM}

Modern visual--inertial estimators are commonly organized either as filters or
as nonlinear optimizers. MSCKF~\cite{mourikis2007msckf} made real-time
filter-based estimation practical by marginalizing feature positions, while
optimization-based approaches jointly refine motion and landmark variables over
a sliding window. OKVIS~\cite{leutenegger2015okvis} demonstrated tightly coupled,
keyframe-based nonlinear optimization for multiple cameras and an IMU.
VINS-Mono~\cite{qin2018vinsmono} added robust initialization, relocalization,
and loop closure in a monocular system, and ORB-SLAM3~\cite{campos2021orbslam3}
unified visual, visual--inertial, and multi-map SLAM. OKVIS2
~\cite{boche2022okvis2} combines a bounded real-time optimization window with
loop-closure constraints and background graph optimization. These systems
provide accurate local motion and can correct drift through revisiting mapped
areas, but loop closure alone does not supply an absolute geographic frame and
cannot be assumed on exploratory trajectories.

\subsection{GNSS-aided visual--inertial estimation}

GNSS and visual--inertial sensing have complementary error characteristics and
have therefore been fused at several levels of coupling. VINS-Fusion
~\cite{qin2019vinsfusion} extends optimization-based VIO to multiple sensor
configurations and global measurements. GVINS~\cite{cao2022gvins} tightly
couples raw GNSS observations with visual and inertial factors, improving the
ability to exploit satellite measurements beyond loose position-fix fusion.
MINS~\cite{geneva2020mins} instead develops an efficient filtering framework
for asynchronous combinations of cameras, IMU, GNSS, wheel encoders, and other
aiding sensors. These representative systems motivate the MINS and
VINS-Fusion baselines in our evaluation, but neither baseline uses a DEM as an
explicit terrain-height constraint in the configuration evaluated here.

Robustness to GNSS failure requires more than adding global position residuals.
Outliers, dropouts, and delayed recovery can invalidate the nominal Gaussian
measurement model. Switchable constraints~\cite{sunderhauf2012switchable} and
dynamic covariance scaling~\cite{agarwal2013dcs} are general robust graph
optimization strategies, while tightly coupled, dropout-tolerant GPS fusion was
developed specifically for visual--inertial SLAM in
~\cite{boche2022dropoutgps}. OKVIS2-X~\cite{boche2025okvis2x} generalizes the
OKVIS2 architecture to optional GNSS and dense depth or LiDAR sensing. Our work
starts from this estimator, but focuses on a different source of global
information: terrain elevation used to stabilize the vertical channel when
GNSS is inaccurate or intermittent. The comparison to default GNSS-enabled
OKVIS2-X is therefore essential for distinguishing the effect of terrain aiding
from that of the underlying estimator.

\subsection{Terrain-aided navigation and elevation priors}

Terrain-aided navigation predates modern visual SLAM. Classical methods compare
measured terrain clearance or elevation profiles with a stored map and estimate
position through nonlinear filtering; Hostetler and Andreas
~\cite{hostetler1983terrain} provided an early treatment of nonlinear Kalman
filtering for this problem. Later work has used Bayesian and particle-filter
formulations to address the nonlinearity and ambiguity of terrain matching
~\cite{bergman1999terrain}. Large-area elevation products such as the Shuttle
Radar Topography Mission~\cite{farr2007srtm} have made globally referenced
terrain data broadly available.

The sensing assumptions of classical terrain-referenced navigation differ from
ours. Many terrain-matching systems use a radar or laser altimeter to observe
clearance directly and use distinctive relief to infer horizontal position. In
contrast, \method{} does not claim to recover horizontal position from a
one-dimensional height residual. It queries the DEM using the available global
horizontal estimate and constrains the platform's expected height after
accounting for a nominal above-ground offset. This design makes the cue usable
inside an existing keyframe graph and prevents the DEM from directly distorting
horizontal visual--inertial motion. It also exposes the method's limitations:
errors in horizontal registration can become height errors on steep slopes,
the nominal clearance assumption is inappropriate for unconstrained flight,
and vegetation, buildings, DEM resolution, and ellipsoidal-versus-orthometric
height conventions can introduce systematic offsets. Our experiments therefore
report performance as a function of slope and GNSS quality and include a
no-DEM ablation rather than treating the map as ground truth.
